\documentclass[10pt]{article}
\usepackage[T1]{fontenc}
\usepackage[utf8]{inputenc}
\usepackage[french]{babel}
\usepackage{amsmath,amssymb,mathrsfs,mathpazo}
\usepackage{geometry}
\geometry{a4paper,margin=2cm}
\usepackage{graphicx,array,tabularx,tabularray,colortbl,multirow,multicol,booktabs,enumitem,pifont,ulem,fancybox,diagbox,cancel}
\usepackage{pstricks,pstricks-add,pst-plot,pst-tree,pst-node,pst-text,pst-func,pst-math,pst-bspline,pst-3dplot,pst-eucl,pst-solides3d}
\let\pscalculate\undefined
\usepackage{pst-bezier}
\usepackage{tikz,pgfplots,tkz-tab}
\usetikzlibrary{arrows,arrows.meta,patterns,calc,positioning,trees,decorations.markings,shapes.geometric,shapes.misc,shapes.arrows,intersections,decorations.pathmorphing,decorations.pathreplacing,backgrounds,shadows,babel}
\pgfplotsset{compat=1.18}
\usepgfplotslibrary{fillbetween}
\usepackage[np]{numprint}
\usepackage{listings,makecell,pas-tableur,twoopt}
\let\barre\undefined
\usepackage[pstricks]{tablvar}
\definecolor{darkgreen}{RGB}{0,100,0}
\providecommand{\red}{\color{red}}
\providecommand{\blue}{\color{blue}}
\providecommand{\green}{\color{green}}
\providecommand{\black}{\color{black}}
\providecommand{\white}{\color{white}}
\providecommand{\gray}{\color{gray}}
\providecommand{\cyan}{\color{cyan}}
\providecommand{\magenta}{\color{magenta}}
\providecommand{\yellow}{\color{yellow}}
\providecommand{\violet}{\color{violet}}
\providecommand{\orange}{\color{orange}}
\providecommand{\euro}{\texteuro}
\providecommand{\EUR}{\texteuro}
\providecommand{\ieme}{\textsuperscript{e}}
\providecommand{\degres}{\ensuremath{^\circ}}
\providecommand{\degree}{\ensuremath{^\circ}}
\providecommand{\vv}[1]{\overrightarrow{#1}}
\providecommand{\Coord}[1]{#1}
\providecommand{\decosix}{}
\providecommand{\starredbullet}{\ensuremath{\bullet}}
\pagestyle{empty}
\setlength{\parindent}{0pt}
\usepackage{adjustbox}\DeclareUnicodeCharacter{2236}{:}\DeclareUnicodeCharacter{2190}{\ensuremath{\leftarrow}}
\providecommand{\R}{\mathbb{R}}
\providecommand{\N}{\mathbb{N}}
\providecommand{\Z}{\mathbb{Z}}
\providecommand{\Q}{\mathbb{Q}}
\providecommand{\D}{\mathbb{D}}
\providecommand{\C}{\mathbb{C}}
\providecommand{\e}{{\mathrm{e}}}
\providecommand{\vect}[1]{\overrightarrow{#1}}
\providecommand{\vectt}[1]{\overrightarrow{\mathrm{#1}}}
\providecommand{\vvt}[1]{\overrightarrow{\mathrm{#1}}}
\providecommand{\np}[2][]{{#2~#1}}
\providecommand{\vv}[1]{{\overrightarrow{#1}}}
\renewcommand{\d}{{\mathrm{d}}}
\providecommand{\degres}{{^\circ}}
\providecommand{\degre}{{^\circ}}
\providecommand{\blue}{\color{blue}}
\providecommand{\red}{\color{red}}
\providecommand{\enskip}{\quad}
\providecommand{\llbracket}{\lbrack\!\lbrack}
\providecommand{\rrbracket}{\rbrack\!\rbrack}
\providecommand{\euro}{{\text{€}}}
\providecommand{\no}{{\text{n°}}}
\providecommand{\interval}[2]{{[#1;#2]}}
\providecommand{\intervFO}[2]{{[#1;#2[}}
\providecommand{\dotfill}{\ldots}

\providecommand{\ds}{\displaystyle}
\providecommand{\textding}[1]{#1}
\providecommand{\Oij}{\(\mathrm{O};\vec{\imath},\vec{\jmath}\)}
\providecommand{\Oijk}{\(\mathrm{O};\vec{\imath},\vec{\jmath},\vec{k}\)}
\providecommand{\Ouv}{\(\mathrm{O};\vec{u},\vec{v}\)}
\providecommand{\intervFF}[2]{\ensuremath{[#1;#2]}}
\renewcommand{\euro}{\ifmmode\text{\texteuro}\else\texteuro\fi}
\DeclareUnicodeCharacter{00B0}{\ifmmode^{\circ}\else\textdegree\fi}
\DeclareUnicodeCharacter{20AC}{\ifmmode\text{\texteuro}\else\texteuro\fi}
% Énoncés : Éducation nationale. Transcription et corrigés : APMEP, Denis Vergès et contributeurs. Découpage et classement : CoopMaths, Sébastien Lozano, Éric Elter et contributeurs. Adaptation de la présentation : Maths Explorer.
% CC BY-SA : https://coopmaths.fr/www/about/
% https://coopmaths.fr/alea/static/bac/2022/tex/bac_2022_05_sujet1_ameriquenord_4_cor.tex
\begin{document}
\noindent\textit{APMEP / CoopMaths. CC BY-SA. Présentation adaptée par Maths Explorer.}\par\medskip
\bigskip
\begin{enumerate}
	\item $\forall x \in \R$, $1-\dfrac{1-\e^x}{1+\e^{x}}=\dfrac{1+\e^x}{1+\e^x}-\dfrac{1-\e^x}{1+\e^x}=\dfrac{2\e^x}{1+\e^x}=\dfrac{2\e^x}{\e^x\left(1+\dfrac{1}{\e^x}\right)}=\dfrac{2}{1+\dfrac{1}{\e^x}}=\dfrac{2}{1+\e^{-x}}$

\hfill \fbox{\textbf{Affirmation 1 : Vraie}}
\item $g(x)=\dfrac12 \iff \dfrac{\e^x}{\e^x+1}=\dfrac12 \iff 2\e^x=\e^x+1 \iff \e^x=1 \iff x=0$
\hfill \fbox{\textbf{Affirmation 2 : Vraie}}
	\item La fonction $f$ est continue et dérivable sur $\R$.

	$\forall x \in \R$, $f'(x)=2x\e^{-x}+x^2\times -\e^{-x}= \left(2x-x^2\right)\e^{-x}$.

	L'équation d'une tangente à $\mathcal{C}$ au point d'abscisse $a$ a pour équation $ y=f'(a)(x-a)+f(a)$. Pour que l'axe des abscisses soit tangent à $\mathcal{C}$, il faut que la courbe $\mathcal{C}$ admette une tangente d'équation $y=0$ (équation de l'axe des abscisses), donc il faut que $f(a)=0$ et $f'(a)=0$. Sachant que $\forall a \in \R$, $\e^{-a}\neq 0$,

$f(a)=0 \iff a^2\e^{-a}=0 \iff a^2=0 \iff a=0$. 

$f'(a)=0\iff (2a-a^2)\e^{-a}=0 \iff 2a-a^2=0 \iff a(2-a)=0 \iff a=0 \text{ ou } a=2$.
	
Donc $a=0$. Il n'existe donc qu'un seul point où l'axe des abscisses est tangent à $\mathcal{C}$.

\hfill \fbox{\textbf{Affirmation 3 : Vraie}}
	\item La fonction $h$ est continue et dérivable sur $\R$. Soit $\mathcal{C}_h$ sa courbe représentative.

$\forall x \in \R$, $h'(x)=\e^x\times \left(1-x^2\right)+\e^x\times-2x=\e^x\left(1-x^2-2x\right)= \left(-x^2-2x+1\right)~\e^x$.

La fonction $h'$ est continue et dérivable sur $\R$.

$\forall x \in \R$, $h''(x)=(-2x-2)\e^x+\e^x\times\left(-x^2-2x+1\right) =\e^x\left(-2x-2-x^2-2x+1\right)= \left(-x^2-4x-1\right)\e^x$

Si $\mathcal{C}_h$ admet des points d'inflexions, alors $h''(x)$ peut s'annuler et changer de signe. Or $\e^x > 0$ pour tout réel $x$, donc étudions le signe de $- x^2- 4x- 1$ sur $\R$.

$-x^2-4x-1$ s'annule pour $x_1=-2-\sqrt{3}$ et pour $x_2=-2+\sqrt{3}$ car $\Delta = 12$.

Le trinôme du second degré s'annule donc deux fois et change donc aussi de signe.

\hfill \fbox{\textbf{Affirmation 4 : Fausse}}
	\item $\forall x \in \R$, $f(x)=\dfrac{\e^x}{\e^x+x}=\dfrac{\e^x}{\e^x\left(1+\dfrac{x}{\e^x}\right)}=\dfrac{1}{1+\dfrac{x}{\e^x}}$.

D'après les croissances comparées, $\lim \limits_{n \rightarrow +\infty} \dfrac{\e^x}{x}=+\infty$ donc $\lim \limits_{n \rightarrow +\infty} \dfrac{x}{\e^x}=0$ donc $\lim \limits_{n \rightarrow +\infty} f(x)=1$

\hfill \fbox{\textbf{Affirmation 5 : Fausse}}
		\item $1+\e^{2x}\geqslant 2\e^x \iff \e^{2x}-2\e^x+1\geqslant0 \iff \left(\e^x-1\right)^2\geqslant0$.

Cette dernière inégalité est vraie pour tout réel $x$.
\hfill \fbox{\textbf{Affirmation 6 : Vraie}}
\end{enumerate}

\end{document}
